% Bewertungspaket zum Gesamttest «Kreis und Kreisteile» — Quelle des PDFs
% (python3 scripts/build-lp-pdf.py kreis-kreisteile). Ganze Punkte je Teilschritt; Werte mit python3 nachgerechnet
% (scripts/lp/kreis-kreisteile/zahlen.py, 08.10.2026, Folgefehler-Fälle eingeschlossen). Aufbau wie
% trigonometrische-berechnungen/bewertungspaket.tex.
\documentclass[11pt]{article}
\usepackage{../lp-druck}
\renewcommand{\lpname}{Kreis und Kreisteile}
\renewcommand{\lpbereich}{Grundlagenfach}
\renewcommand{\lpteilgebiet}{5.2c}
\renewcommand{\lpfuss}{Bewertungspaket}
\newcolumntype{P}{>{\centering\arraybackslash}p{10mm}}
\newenvironment{raster}{\par\smallskip\noindent\tabularx{\linewidth}{@{}>{\raggedright\arraybackslash}p{38mm}P X@{}}\toprule
  \small\textsc{Teilschritt} & \small\textsc{P} & \small\textsc{Musterlösung}\\\midrule}{\bottomrule\endtabularx\par}
\newcommand{\fehler}[1]{\par\smallskip{\small\textcolor{lprot}{\bfseries Typische Fehler:} #1}}
\newcommand{\E}{\,{\footnotesize\bfseries\color{lpgruen}(E)}}
\newcommand{\A}{\,{\footnotesize\bfseries\color{lpblau}(A)}}
\begin{document}
\lpkopf{Bewertungspaket zum Gesamttest}{}

\begin{lpachtung}{Erst lösen, dann öffnen}
Dieses Paket enthält die Musterlösung. Wer es vor dem Lösen liest, misst am Ende nur, wie gut das Abschreiben gelingt.
\end{lpachtung}

\section*{1 · So gehst du vor}
\begin{enumerate}[leftmargin=*,itemsep=0pt]
\item \textbf{Gesamttest lösen} — vollständig, mit Skizze und Rechenweg, ohne dieses Paket; der Taschenrechner ist erlaubt.
\item \textbf{Fotografieren oder scannen}, gut lesbar, eine Seite pro Bild. \textbf{Kein Name} auf den Bildern.
  Handyfotos tragen oft unsichtbar Standort, Zeit und Gerät mit (Metadaten). Lade darum lieber einen Scan oder ein
  Bildschirmfoto des Fotos hoch, oder entferne vorher die Standortangaben.
\item \textbf{Bewerten lassen.} Ob du dafür eine KI nutzt und welche, entscheidest du selbst und in eigener Verantwortung —
  je nachdem, was dir aufgrund deines Alters und deiner persönlichen Situation erlaubt ist (Altersgrenzen und
  Nutzungsbedingungen des Anbieters, allenfalls Einverständnis deiner Eltern). Lade nur deine Lösung und dieses PDF hoch,
  keine persönlichen Daten, und füge den Auftrag aus Abschnitt~2 ein. Ohne KI kannst du dich mit Abschnitt~3 selbst bewerten.
\item \textbf{Rückmeldung prüfen.} Stimmt, was die KI aus deiner Handschrift gelesen hat? Eine KI kann sich irren —
  im Zweifel gilt das Raster in Abschnitt~3.
\item \textbf{Gezielt wiederholen} nach Abschnitt~4.
\end{enumerate}

\needspace{8\baselineskip}
\section*{2 · Auftrag an die KI}
\begin{tcolorbox}[breakable,colback=lplinie!25,colframe=lplinie,boxrule=0.5pt,arc=1mm,fontupper=\small,left=2mm,right=2mm,top=1.5mm,bottom=1.5mm]
Du bist Korrektorin oder Korrektor für Mathematik an einer Schweizer Berufsmaturitätsschule. Im Anhang findest du (1)~das Bewertungspaket zum Gesamttest «Kreis und Kreisteile» mit Musterlösung und Punkteraster und (2)~Fotos meiner handschriftlichen Lösung.\par\smallskip
Geh so vor:\par
1. Schreib zu jeder Aufgabe G1 bis G7 zuerst ab, was du in meiner Lösung liest — Skizze, Rechenweg und Ergebnis. Was du nicht sicher lesen kannst, markierst du mit [unsicher]. Nicht raten.\par
2. Vergib die Punkte genau nach dem Raster im Bewertungspaket (Abschnitt 3), ganze Punkte je Zeile. Ziehe einen Fehler nur einmal ab: Rechne ich mit einem falschen Zwischenergebnis richtig weiter, gibt es die Folgepunkte. Ausnahme: Zeilen mit \textbf{(E)} gibt es nur für das richtige Ergebnis, nie als Folgepunkt. Die Zeilen «Typische Fehler» sagen, welche Rasterzeile bei diesem Fehler 0 Punkte gibt — sie sind kein zusätzlicher Abzug.\par
3. Andere richtige Lösungswege zählen voll, ebenso gleichwertige Schreibweisen: Dezimalpunkt oder Dezimalkomma, exakte Werte wie $6\pi$ oder $\sqrt{50}$ statt Dezimalzahlen, anders, aber vernünftig gerundete Ergebnisse ($\pm 0.01$ oder durch gerundete Zwischenwerte erklärbar). \textbf{Einheiten:} Fehlt bei einem Endergebnis die Einheit oder ist sie falsch ($\text{cm}$ statt $\text{cm}^2$), gibt die (E)-Zeile dieses Ergebnisses 0 Punkte — aber nur beim ersten Mal im ganzen Gesamttest; weitere Einheitenfehler kosten nichts mehr. Zeilen mit \textbf{(A)} verlangen nur Erkennen, diese Punkte gibt es auch ohne Rechenweg. Zeilen mit \textbf{(E)} sind Ergebnispunkte: Das Ergebnis muss stimmen \emph{und} aus einem erkennbaren Weg hervorgehen. Alle übrigen Zeilen bewerten den Weg selbst.\par
4. Begründe jeden Abzug in einem Satz und nenne die Fehlerart (zum Beispiel «Durchmesser statt Radius» oder «Dreieck abgezogen statt addiert»). Schreib die Musterlösung nicht ab — zeig mir, wo mein Fehler liegt.\par
5. Gib zum Schluss eine Tabelle «Aufgabe | Punkte | Begründung», die Gesamtpunktzahl (von 25) und — nach der Selbsteinschätzung in Abschnitt 4 — welches Kapitel des Leitprogramms ich wiederholen soll. Keine Note.\par
6. Fehlt ein Foto oder ist eine Aufgabe unlesbar, sag es, statt sie zu bewerten.
\end{tcolorbox}

\needspace{8\baselineskip}
\section*{3 · Musterlösung und Punkteraster}
\textbf{Allgemein:} Ganze Punkte je Zeile. Ein Fehler wird nur einmal abgezogen (Folgefehler: wer mit einem falschen
Zwischenergebnis richtig weiterrechnet, bekommt die Wegpunkte der folgenden Zeilen, aber keine (E)-Punkte). In der Spalte P:
\textbf{(A)}~= erkennen, zählt auch ohne Rechenweg · \textbf{(E)}~= Ergebnispunkt, Ergebnis richtig
\emph{und} Weg erkennbar · ohne Zeichen = die Zeile bewertet den Weg. Taschenrechner erlaubt; Ergebnisse auf zwei
Dezimalen, Toleranz $\pm 0.01$ (oder durch gerundete Zwischenwerte erklärbar). \textbf{Einheiten:} Der erste Einheitenfehler im
ganzen Test kostet die (E)-Zeile dieses Ergebnisses, weitere nichts. \textbf{$\pi$:} Gerechnet wird mit der $\pi$-Taste. Mit
$\pi \approx 3.14$ zählt der Weg; die (E)-Zeile nur, wenn das Ergebnis trotzdem auf $\pm 0.01$ stimmt.
\textbf{Kein Folgepunkt für Unmögliches:} Ein Zentriwinkel über $360^\circ$ oder eine Sehne, die länger ist als der Durchmesser,
gibt keine Folgepunkte. «Typische Fehler» nennt, welche Zeile 0 Punkte gibt, und die Punkte, die bleiben.
\textbf{Teilrichtige Zeilen:} Verlangt eine Zeile mehrere Angaben und stimmt nur ein Teil, gibt die Zeile 0 Punkte.

\begin{lpaufgabe}{G1}{Gerade und Kreis}{3 P}
\begin{raster}
(a) Lage & 1\A & $r = \tfrac{d}{2} = 4\,\text{cm}$. $g$: $a_g = r$, Tangente. $h$: $a_h < r$, Sekante. $k$: $a_k > r$, Passante. Alle drei mit dem Vergleich zu $r$ nötig.\\
(b) Ansatz & 1 & Das Lot von $M$ halbiert die Sehne: rechtwinkliges Dreieck mit der Hypotenuse $r$, $\left(\tfrac{s}{2}\right)^2 + 2.5^2 = 4^2$. Folgepunkt: richtig mit dem eigenen Radius.\\
(b) Sehne & 1\E & $\tfrac{s}{2} = \sqrt{16 - 6.25} = \sqrt{9.75} \approx 3.12\,\text{cm}$, $s \approx 6.24\,\text{cm}$\\
\end{raster}
\fehler{Durchmesser als Radius ($r = 8$): (a) dann alle drei «Sekante» — 0; (b) $2\sqrt{64 - 6.25} \approx 15.20$: Ansatz als Folgepunkt, Ergebnis 0 $\to$ 1 von 3\,P.
Nur die halbe Sehne $3.12\,\text{cm}$ als Ergebnis: (b)-Ergebniszeile 0 $\to$ 2 von 3\,P.
$\sqrt{4^2 + 2.5^2}$ (plus statt minus): (b) Ansatz und Ergebnis 0.}
\end{lpaufgabe}

\begin{lpaufgabe}{G2}{Wie weit reicht der Blick?}{3 P}
\begin{raster}
Skizze & 1 & Kreis (Erde) mit Mittelpunkt $M$, die Lampe $P$ ausserhalb auf der Verlängerung eines Radius, der Lichtstrahl als Tangente von $P$ an den Kreis mit dem Berührpunkt $B$, rechter Winkel bei $B$ eingezeichnet.\\
Ansatz & 1 & Rechter Winkel bei $B$ (Tangente $\perp$ Radius), $\overline{MP} = 6371\,\text{km} + 0.05\,\text{km} = 6371.05\,\text{km}$ ist die Hypotenuse: $\overline{PB}^2 = \overline{MP}^2 - r^2$. Beides nötig: Hypotenuse $\overline{MP}$ und die Einheiten angeglichen.\\
Ergebnis & 1\E & $\overline{PB} = \sqrt{6371.05^2 - 6371^2} \approx 25.24\,\text{km}$\\
\end{raster}
\fehler{$50\,\text{m}$ nicht umgerechnet ($\overline{MP} = 6421$): $\sqrt{6421^2 - 6371^2} \approx 799.75\,\text{km}$ — Ansatz und Ergebnis 0, die Skizze zählt $\to$ 1 von 3\,P.
$\sqrt{\overline{MP}^2 + r^2} \approx 9009.99\,\text{km}$ ($\overline{MP}$ als Kathete): Ansatz und Ergebnis 0.
$0.05\,\text{km}$ als Antwort (Höhe statt Sichtweite): Ansatz und Ergebnis 0.}
\end{lpaufgabe}

\begin{lpaufgabe}{G3}{Teich mit Weg}{4 P}
\begin{raster}
(a) Radius & 1 & $U = 2\pi r$: $r = \tfrac{25}{2\pi} \approx 3.98\,\text{m}$\\
(a) Fläche & 1\E & $A = \pi r^2 \approx 49.74\,\text{m}^2$ (mit $r = 3.98$ gerechnet $49.76\,\text{m}^2$: zählt)\\
(b) Ansatz & 1 & Der Weg ist ein Kreisring mit $r \approx 3.98\,\text{m}$ und $R = r + 1.5 \approx 5.48\,\text{m}$: $A = \pi (R^2 - r^2)$ — gleichwertig $2\pi r_m \cdot b$ mit $r_m \approx 4.73\,\text{m}$, $b = 1.5\,\text{m}$. Folgepunkt: richtig mit dem eigenen Radius.\\
(b) Weg & 1\E & $A \approx 44.57\,\text{m}^2$ (mit gerundeten Radien $44.58\,\text{m}^2$: zählt)\\
\end{raster}
\fehler{$r = \tfrac{25}{\pi} \approx 7.96$ (das ist der Durchmesser): Radiuszeile 0, (a) Fläche $198.94$ 0; (b) Ansatz als Folgepunkt ($82.07$), Ergebnis 0 $\to$ 1 von 4\,P.
Weg als $\pi \cdot 1.5^2 \approx 7.07$: beide (b)-Zeilen 0.
$\pi (R - r)^2$: beide (b)-Zeilen 0.}
\end{lpaufgabe}

\begin{lpaufgabe}{G4}{Ein Sektor, rückwärts}{4 P}
\begin{raster}
(a) Ansatz & 1 & Anteil $= \tfrac{A_{SK}}{\pi r^2} = \tfrac{50}{\pi \cdot 7.5^2}$, dann $\varphi = \text{Anteil} \cdot 360^\circ$ (gleichwertig: $50 = \tfrac{\varphi}{360^\circ} \cdot \pi \cdot 7.5^2$ nach $\varphi$ aufgelöst)\\
(a) Winkel & 1\E & $\varphi \approx 101.86^\circ$\\
(b) Bogen & 1 & $b = \tfrac{2 A_{SK}}{r} = \tfrac{100}{7.5} \approx 13.33\,\text{cm}$ (aus $A_{SK} = \tfrac12 b \cdot r$) oder $b = \tfrac{\varphi}{360^\circ} \cdot 2\pi r$. Folgepunkt: richtig mit dem eigenen $\varphi$, sofern $\varphi < 360^\circ$.\\
(c) Rand & 1 & $U_S = b + 2r \approx 13.33 + 15 = 28.33\,\text{cm}$. Folgepunkt: eigenes $b$ plus zwei Radien.\\
\end{raster}
\fehler{Anteil mit dem Umfang ($\tfrac{50}{2\pi \cdot 7.5} \cdot 360^\circ \approx 381.97^\circ$): (a) beide Zeilen 0; ein Winkel über $360^\circ$ gibt keine Folgepunkte, (b) zählt nur, wenn $b$ über $\tfrac{2A_{SK}}{r}$ richtig berechnet ist $\to$ (c) als Folgepunkt.
Anteil $0.28$ als Winkel stehen gelassen: (a) Winkelzeile 0.
Rand ohne die Radien ($13.33$) oder mit nur einem ($20.83$): (c) 0.}
\end{lpaufgabe}

\begin{lpaufgabe}{G5}{Ein Segment}{4 P}
\begin{raster}
(a) Radius & 1 & Das Lot halbiert die Sehne: rechtwinkliges Dreieck mit den Katheten $5$ (halbe Sehne) und $5$ (Abstand): $r = \sqrt{5^2 + 5^2} = \sqrt{50} \approx 7.07\,\text{cm}$\\
(b) Begründung & 1 & Das halbe Dreieck ist rechtwinklig mit zwei gleich langen Katheten, also gleichschenklig: Sein Winkel bei $M$ ist $45^\circ$, die beiden Hälften zusammen $90^\circ$. Gleichwertig: $r^2 + r^2 = 50 + 50 = 100 = s^2$, nach der Umkehrung des Pythagoras liegt bei $M$ ein rechter Winkel.\\
(c) Sektor & 1 & $A_{SK} = \tfrac{90^\circ}{360^\circ} \cdot \pi \cdot 50 = 12.5\pi \approx 39.27\,\text{cm}^2$. Folgepunkt: richtig mit dem eigenen $r$.\\
(c) Segment & 1\E & Dreieck $\tfrac12 \cdot r \cdot r = 25\,\text{cm}^2$ (gleichwertig $\tfrac12 \cdot 10 \cdot 5$); $A_{SG} = 39.27 - 25 \approx 14.27\,\text{cm}^2$\\
\end{raster}
\fehler{Den Abstand als Radius genommen ($r = 5$): (a) 0; (c) Sektor als Folgepunkt ($19.63$), Segment 0 $\to$ höchstens 2 von 4\,P.
Sektor plus Dreieck ($64.27$): Segmentzeile 0.
(b) nur «sieht rechtwinklig aus» oder ein Messen mit dem Geodreieck: 0.}
\end{lpaufgabe}

\begin{lpaufgabe}{G6}{Eine fremde Aussage}{3 P}
\begin{raster}
Urteil & 1 & Nein: Die Fläche wird \emph{achtmal} so gross. Nötig: «nein» \emph{und} der Faktor $8$.\\
Begründung & 1 & $A_{SK} = \tfrac{\varphi}{360^\circ} \cdot \pi r^2$: doppelter Winkel gibt den Faktor $2$, doppelter Radius den Faktor $2^2 = 4$, zusammen $2 \cdot 4 = 8$. Ein Zahlenbeispiel zählt, wenn es beide Verdopplungen zeigt.\\
Bogen & 1 & $b = \tfrac{\varphi}{360^\circ} \cdot 2\pi r$: Faktor $2 \cdot 2 = 4$ — der Bogen wird viermal so lang.\\
\end{raster}
\fehler{«Stimmt, Faktor 4» (nur der Radius berücksichtigt): Urteil und Begründung 0; die Bogenzeile nach ihrer Zeile.
«Faktor 6» ($2 + 4$): Urteil und Begründung 0.}
\end{lpaufgabe}

\begin{lpaufgabe}{G7}{Sportplatz}{4 P}
\begin{raster}
(a) Ansatz & 1 & Rand: die beiden langen Seiten und zwei Halbkreisbögen, die zusammen einen ganzen Kreisumfang mit $d = 64\,\text{m}$ ergeben: $2 \cdot 100 + \pi \cdot 64$\\
(a) Rand & 1\E & $\approx 401.06\,\text{m}$\\
(b) Ansatz & 1 & Rechteck und zwei Halbkreise (zusammen ein Kreis mit $r = 32\,\text{m}$): $100 \cdot 64 + \pi \cdot 32^2$\\
(b) Fläche & 1\E & $\approx 9616.99\,\text{m}^2$\\
\end{raster}
\fehler{Die schmalen Rechteckseiten mitgezählt ($529.06\,\text{m}$): (a) Ansatz und Ergebnis 0.
$\pi d^2$ statt $\pi r^2$ ($19\,267.96\,\text{m}^2$) oder nur ein Halbkreis ($8008.50\,\text{m}^2$; beim Rand $300.53\,\text{m}$): die betroffene Ansatz- und Ergebniszeile 0.
Mit $\pi \approx 3.14$ ($400.96\,\text{m}$; $9615.36\,\text{m}^2$): Ansatzzeilen zählen, (E)-Zeilen 0 $\to$ 2 von 4\,P.}
\end{lpaufgabe}

\needspace{14\baselineskip}
\section*{4 · Selbsteinschätzung}
\begin{tabularx}{\linewidth}{@{}l X@{}}\toprule
22 – 25 P & Die geprüften Teile sitzen. Wo du Punkte verloren hast: das Kapitel dieser Aufgabe nochmals (Zuordnung unten).\\
17 – 21 P & Das schwächste Kapitel nochmals: Tüfteln und Übungen des Kapitels, in dem du die meisten Punkte verloren hast.\\
11 – 16 P & Zurück zu den Kapiteln aller Aufgaben, in denen du Punkte verloren hast.\\
0 – 10 P & Zurück zu Kapitel 1 und von dort der Reihe nach weiter.\\\midrule
\multicolumn{2}{@{}p{\linewidth}@{}}{\small\textbf{Aufgabe $\to$ Kapitel:} G1, G2 $\to$ Kapitel 1 (Linien am Kreis); G3 (a) $\to$ Kapitel 2 (Umfang, Fläche und $\pi$), G3 (b) $\to$ Kapitel 4 (Kreisring); G4, G6 $\to$ Kapitel 3 (Bogen und Sektor); G5 (a) $\to$ Kapitel 1, G5 (b), (c) $\to$ Kapitel 4 (Segment); G7 $\to$ Kapitel 2 und 4 (zusammengesetzte Flächen)}\\\bottomrule
\end{tabularx}
\end{document}
